% ID: SEL-D0014-BD-E07
% ESTADO: validada. S015; tres pasadas del mismo autor, sin segundo revisor.
\textbf{a) Distribución de las medias.}\par\nopagebreak[4]
La varianza poblacional es $9$, luego $\sigma=3$ euros. Para observaciones independientes de una población normal, la media muestral de tamaño $10$ es normal exacta, con
\[
\boxed{E(\overline X)=\mu,\qquad
\operatorname{Var}(\overline X)=\frac9{10},\qquad
\sigma_{\overline X}=\frac3{\sqrt{10}}.}
\]
Se explicitan varianza y desviación para evitar la ambigüedad de la segunda coordenada de la notación normal.

\textbf{b) Intervalo para el precio medio.}\par\nopagebreak[4]
Los diez precios suman $110$ euros, así que $\overline x=11$ euros. Para confianza 97\,\%, $z_{0.985}\approx2.170090378$. Al ser conocida la desviación poblacional,
\[
IC=11\pm z_{0.985}\frac3{\sqrt{10}},\qquad E\approx2.058728496\text{ euros}.
\]
\[
\boxed{IC_{\text{97\%}}\approx(8.941271504,\ 13.058728496)\text{ euros/kg}.}
\]
La confianza describe la cobertura del procedimiento bajo el modelo de muestreo, no una probabilidad posterior de la media fija.

\textbf{c) Error estrictamente menor que 1.5 euros.}\par\nopagebreak[4]
Se conservan el cuantil y $\sigma$:
\[
z_{0.985}\frac3{\sqrt n}<1.5
\quad\Longleftrightarrow\quad
n>\left(\frac{3z_{0.985}}{1.5}\right)^2
\approx18.837168988.
\]
\boxed{n_{\min}=19}. Para $18$ el error no es menor que $1.5$ y para $19$ sí. En general, la desigualdad estricta exige $\lfloor\text{cota}\rfloor+1$, incluso si ésta fuese entera. Con $2.17$ tabulado se mantiene $19$.
